\section{Model, transport restrictions and inference criterion}
\label{sec:model}

Throughout, sample sizes are positive integers, and all random variables and conditional expectations are measurable. We use \(\lVert\cdot\rVert_\infty\) for the supremum norm and \(|\cdot|\) for absolute value. Generic positive constants may depend on the fixed density envelopes, smoothness radius and coverage level; their values may change between occurrences. The notation \(u_n\lesssim v_n\) denotes an inequality up to such a constant, uniformly over the stated model class and strength range, and \(u_n\asymp v_n\) denotes inequalities in both directions.

\subsection{Primitive variables and the observed experiment}

The scalar covariate domain is \leanref{sym:\mathcal X}{\ensuremath{\mathcal X}}=[0,1]. Population membership is coded by \(S=1\) for the source and \(S=0\) for the target. A primitive full-data record under the full-data law \leanref{sym:P_n^F}{\ensuremath{P_n^F}} is
\[
O^F=(S,X,D(0),D(1),Y(0),Y(1)),
\]
where \(D(0),D(1)\in\{0,1\}\) are potential receipts under encouragement, \(Y(0),Y(1)\in[0,1]\) are potential outcomes under receipt, and \(C=\mathbf 1\{D(1)>D(0)\}\) is the complier indicator. The target conditional effect among compliers is represented by
\[
\theta=
\frac{\mathbb E_{P_n^F}[(Y(1)-Y(0))C\mid S=0]}
     {\mathbb E_{P_n^F}[C\mid S=0]},
\]
whenever the denominator is positive.

The source assignment law augments a source full-data draw with encouragement \(Z\), using propensity \leanref{sym:e}{\ensuremath{e}}, and consistency induces the observed receipt \(D=D(Z)\) and outcome \(Y=Y(D)\). Its observed record is \(O^{\mathrm S}=(X,Z,D,Y)\) with one-record law \leanref{sym:P_{\mathrm S}}{\ensuremath{P_{\mathrm S}}}. The target observation is only \(X^{\mathrm T}=X\), with one-record law \leanref{sym:P_{\mathrm T}}{\ensuremath{P_{\mathrm T}}}. Thus \(P_n^F\) determines the causal variables, the source assignment mechanism determines the observed source law, and the target-population covariate marginal determines \(P_{\mathrm T}\). The samples \leanref{sym:\mathcal S_n}{\ensuremath{\mathcal S_n}} and \leanref{sym:\mathcal T_n}{\ensuremath{\mathcal T_n}} contain \(n\) records from these two observed laws. Their densities and conditional arm means receive the regularity restrictions below.

\begin{assumptionv}{ass:source-iid}[Independent source sampling]
The observed source records \(\mathcal S_n\) are independent and identically distributed with common law \(P_{\mathrm S}\), the law of one observed source record:
\[
\mathcal S_n\sim P_{\mathrm S}^{\otimes n}.
\]
\end{assumptionv}

We first impose independent sampling within each population and independence between the two samples.

\begin{assumptionv}{ass:target-iid}[Independent target sampling]
The observed target covariate records \(\mathcal T_n\) are independent and identically distributed with common law \(P_{\mathrm T}\), the law of one observed target covariate:
\[
\mathcal T_n\sim P_{\mathrm T}^{\otimes n}.
\]
\end{assumptionv}

Independent source sampling makes the source records repeated observations from a common population law, as in standard two-sample transport designs \citep{ZengKennedyBodnarNaimi2025Efficient}. The corresponding target condition fixes the population represented by the covariate sample.

\begin{assumptionv}{ass:sample-independence}[Source–target sample independence]
The observed source records \(\mathcal S_n\) and observed target covariate records \(\mathcal T_n\) are independent:
\[
\mathcal S_n\perp\!\!\!\perp\mathcal T_n.
\]
\end{assumptionv}

Independent target sampling allows the target covariate distribution to be learned from repeated draws under the same two-sample framework \citep{ZengKennedyBodnarNaimi2025Efficient}. We complete the sampling design by requiring the two collections of records to be independent.

\begin{definitionv}{synth_1}[The scalar covariate space]
The covariate space is
\[
\mathcal X := [0,1] \subset \mathbb R.
\]
\end{definitionv}

Source–target sample independence describes separately drawn samples and gives the joint observed law a product structure \citep{ZengKennedyBodnarNaimi2025Efficient}.

\subsection{Density, overlap and smoothness restrictions}

The regularity conditions use fixed envelopes and a common Hölder convention. We define these quantities together so that density, propensity and regression restrictions share the same scale.



The lower density envelope \leanref{sym:c_f}{\ensuremath{c_f}} and upper density envelope \leanref{sym:C_f}{\ensuremath{C_f}} control covariate support across populations. The common radius \leanref{sym:L}{\ensuremath{L}} bounds both components of the Hölder ball \leanref{sym:\mathcal H^{\beta}([0,1],L)}{\ensuremath{\mathcal H^\beta([0,1],L)}}. At exponent \(\beta=1/8\), the class permits rough continuous nuisance functions. We first impose density bounds for the source and target designs.

\begin{assumptionv}{ass:source-density-bounds}[Source density and support bounds]
The source covariate density \(f_{\mathrm S}\), with \(c_f\) and \(C_f\) the fixed lower and upper density envelopes, satisfies the following conditions:
\begin{itemize}
\item \textbf{(Envelope constants.)} \(0<c_f<1<C_f\).
\item \textbf{(Full-data law and support.)} \(P_n^F\), the full-data law, is a probability law. Under this law, the covariate \(X\) and both potential outcomes \(Y(0)\) and \(Y(1)\) belong to \([0,1]\) almost surely.
\item \textbf{(Density identity.)} The covariate marginal of \(P_{\mathrm S}\), the law of one observed source record, has density \(f_{\mathrm S}\) with respect to Lebesgue measure restricted to \(\mathcal X=[0,1]\). Thus, for every Borel set \(B\subseteq\mathbb R\),
\[
P_{\mathrm S}(X\in B)=\int_{B\cap[0,1]} f_{\mathrm S}(x)\,dx.
\]
\item \textbf{(Pointwise bounds.)} For every \(x\in\mathcal X\),
\[
c_f\le f_{\mathrm S}(x)\le C_f.
\]
\end{itemize}
\end{assumptionv}

A uniformly positive bounded source design density supplies source observations throughout the covariate domain, a standard support condition for transport analysis \citep{ZengKennedyBodnarNaimi2025Efficient}. The full-data support clause also fixes the scale of the potential outcomes. The target density obeys the same envelopes.

\begin{assumptionv}{ass:target-density-bounds}[Target covariate density bounds]
The target covariate law \(P_{\mathrm T}\), with target density \(f_{\mathrm T}\), satisfies the following conditions on \(\mathcal X=[0,1]\), the covariate domain, with real constants \(c_f\) and \(C_f\), the lower and upper density envelopes:
\begin{itemize}
\item \textbf{(Density identity.)} For every Borel set \(B\subseteq\mathbb R\),
\[
P_{\mathrm T}(B)=\int_{B\cap\mathcal X}\max\{f_{\mathrm T}(x),0\}\,dx.
\]
\item \textbf{(Density envelope.)} For every \(x\in\mathcal X\),
\[
c_f\le f_{\mathrm T}(x)\le C_f.
\]
\item \textbf{(Positive lower envelope.)}
\[
0<c_f.
\]
\end{itemize}
\end{assumptionv}

Uniformly positive bounded target density gives both populations the common support used in transport comparisons \citep{ZengKennedyBodnarNaimi2025Efficient}. Together, the two density envelopes bound the target-to-source density ratio. We next control variation within that support.

\begin{assumptionv}{ass:source-density-holder}[Source density Hölder regularity]
The source covariate density \(f_{\mathrm S}\) belongs to \(\mathcal H^{\beta}([0,1],L)\), the class of continuous functions on \([0,1]\) with supremum norm and \(\beta\)-Hölder seminorm bounded separately by \(L\), where \(\beta\) is the common Hölder exponent and \(L\) is the common Hölder radius:
\[
\beta=\frac18,\qquad L>1,\qquad
f_{\mathrm S}\in\mathcal H^{\beta}([0,1],L).
\]
Thus the bounds are
\[
\lVert f_{\mathrm S}\rVert_{\infty}\le L,
\qquad
[f_{\mathrm S}]_{\beta}\le L,
\]
where \([f_{\mathrm S}]_{\beta}\) denotes the \(\beta\)-Hölder seminorm of the source density.
\end{assumptionv}

The one-dimensional isotropic Hölder ball controls the source density's local variation while allowing low regularity, the smoothness framework used in higher-order functional estimation \citep{RobinsLiMukherjeeTchetgenVanDerVaart2017HigherOrder}. The target density receives the same exponent and radius.

\begin{assumptionv}{ass:target-density-holder}[Target density Hölder regularity]
The target covariate density \(f_{\mathrm T}\) satisfies the following conditions:
\begin{itemize}
\item \textbf{(Radius.)} The common Hölder radius \(L\) satisfies \(L>1\).
\item \textbf{(Exponent.)} The common Hölder exponent is \(\beta=1/8\).
\item \textbf{(Regularity.)} \(f_{\mathrm T}\in\mathcal H^\beta([0,1],L)\): it is continuous on \([0,1]\) and satisfies
\[
\sup_{x\in[0,1]}|f_{\mathrm T}(x)|\le L,
\qquad
[f_{\mathrm T}]_\beta
=
\sup_{\substack{x,y\in[0,1]\\x\ne y}}
\frac{|f_{\mathrm T}(x)-f_{\mathrm T}(y)|}{|x-y|^\beta}
\le L.
\]
\end{itemize}
\end{assumptionv}

Target density Hölder regularity places the target design in the same one-dimensional smoothness class as the source density, giving common control over the covariate distributions entering our transported functional. The related transport rate comparison specifies smoothness of outcome regressions and inverse selection--treatment weights \citep{ZengKennedyBodnarNaimi2025Efficient}. For source encouragement, we impose positivity before its smoothness restriction.

\begin{assumptionv}{ass:propensity-overlap}[Source encouragement propensity overlap]
The source encouragement propensity \(e\) satisfies
\[
\frac14\le e(x)\le\frac34
\qquad\text{for every }x\in\mathcal X.
\]
\end{assumptionv}

Source instrument positivity ensures that both encouragement arms are represented at every covariate value \citep{ChenHuang2025Generalizing}. Its fixed bounds concern assignment probabilities, whereas compliance strength will be governed by a separate condition. The propensity also satisfies the common Hölder restriction.

\begin{assumptionv}{ass:propensity-holder}[Hölder regularity of the propensity]
The source encouragement propensity \(e\) belongs to \(\mathcal H^\beta([0,1],L)\), with Hölder exponent \(\beta=1/8\) and common Hölder radius \(L>1\); explicitly,
\[
\lVert e\rVert_\infty\le L,
\qquad
[e]_\beta\le L.
\]
\end{assumptionv}

A Hölder-smooth propensity controls how assignment probabilities vary with covariates in our model. The related transport rate comparison imposes smoothness on an inverse selection--treatment weight and an outcome regression \citep{ZengKennedyBodnarNaimi2025Efficient}. The remaining observed nuisance functions are the receipt and outcome regressions within each encouragement arm.

\begin{assumptionv}{ass:arm-means-holder}[Hölder regularity of arm means]
The source conditional arm means \(m_{Az}\), for responses \(A\in\{Y,D\}\) and encouragement arms \(z\in\{0,1\}\), satisfy the following conditions, with common Hölder radius \(L\) and Hölder exponent \(\beta\):
\begin{itemize}
\item \textbf{(Smoothness.)} \(L>1\), \(\beta=1/8\), and each \(m_{Az}\) is continuous on \([0,1]\), with
\[
\lVert m_{Az}\rVert_\infty\le L,
\qquad
[m_{Az}]_\beta\le L.
\]
\item \textbf{(Range.)} For every \(A\), \(z\), and \(x\in[0,1]\),
\[
0\le m_{Az}(x)\le 1.
\]
\item \textbf{(Source arm moments.)} Under the usual measurability conditions, for every \(A\), \(z\), and measurable set \(B\subseteq[0,1]\),
\[
\int_{\{X\in B,\;Z=z\}} A\,dP_{\mathrm S}
=
\int_B f_{\mathrm S}(x)\,
e(x)^z\bigl(1-e(x)\bigr)^{1-z}\,
m_{Az}(x)\,dx,
\]
where \(P_{\mathrm S}\) is the law of an observed source record, \(f_{\mathrm S}\) is the source covariate density, and \(e\) is the source encouragement propensity.
\end{itemize}
\end{assumptionv}

Hölder-smooth observed arm regressions control the covariate variation of the conditional receipt and outcome means \citep{ZengKennedyBodnarNaimi2025Efficient}. Their range reflects the bounded responses, and the arm-moment identities tie the chosen regression versions to the observed source law.

\subsection{Instrument validity and complier transport}

The causal restrictions connect observed encouragement-arm contrasts to effects among compliers. We begin with consistency of treatment receipt and outcomes in the source population.

\begin{assumptionv}{ass:receipt-consistency}[Source treatment receipt consistency]
In the source population, observed receipt equals potential receipt under the realized encouragement:
\[
D=D(Z)\qquad\text{almost surely}.
\]
\end{assumptionv}

Receipt consistency connects realized treatment receipt to its potential value under the assigned encouragement \citep{ChenHuang2025Generalizing}. The outcome restriction completes this link between observed records and potential variables.

\begin{assumptionv}{ass:outcome-consistency}[Outcome consistency and boundedness]
In the source population, the observed outcome equals the potential outcome under the realized receipt and lies in the unit interval, both almost surely:
\[
Y=Y(D), \qquad 0\le Y\le 1.
\]
\end{assumptionv}

Outcome consistency identifies the observed response with the potential outcome under realized receipt \citep{ChenHuang2025Generalizing}; boundedness fixes its scale for the inference problem. We then require conditional randomization of encouragement.

\begin{assumptionv}{ass:instrument-randomization}[Conditional instrument randomization]
Let \(O\) denote the primitive full-data record and let \(P_n^{F,\mathrm S}=P_n^F(\,\cdot\mid S=1)\) be its source-population law. The joint source assignment law \(\mathbb P_{\mathrm{assign}}\) of \((O,Z)\) is a probability law, and its \(O\)-marginal is \(P_n^{F,\mathrm S}\). The zero extension of the source propensity,
\[
\widetilde e(x)=
\begin{cases}
e(x),&x\in[0,1],\\
0,&x\notin[0,1],
\end{cases}
\]
is measurable. Put \(\pi_1(x)=e(x)\) and \(\pi_0(x)=1-e(x)\). For every measurable full-data event \(B\) and each \(z\in\{0,1\}\),
\[
\mathbb P_{\mathrm{assign}}(O\in B, Z=z)
=\int_B \pi_z(X(O))\,dP_n^{F,\mathrm S}(O).
\]
Thus encouragement \(Z\) is conditionally independent of \((D(0),D(1),Y(0),Y(1))\) given \(X\) and source membership \(S=1\):
\[
Z\perp\!\!\!\perp (D(0),D(1),Y(0),Y(1))\mid X,S=1.
\]
\end{assumptionv}

Conditional instrument exchangeability makes source encouragement comparable across potential receipt and outcome types within covariate strata \citep{ChenHuang2025Generalizing}. This permits a covariate-dependent randomized encouragement design. The next restriction specifies the causal channel through which encouragement affects outcomes.

\begin{assumptionv}{ass:exclusion}[Exclusion through treatment receipt]
For each \(z\in\{0,1\}\), the source potential outcome generated by setting encouragement to \(z\) equals \(Y(D(z))\) almost surely under \(P_n^F\), the primitive full-data law at sample-size index \(n\). Here \(D(z)\) is potential receipt under encouragement \(z\), and \(Y(d)\) is the potential outcome under receipt \(d\).
\end{assumptionv}

Instrument exclusion assigns encouragement's outcome effect to the receipt it induces \citep{ChenHuang2025Generalizing}. Combined with conditional randomization and consistency, it supplies the outcome interpretation of the source arm contrast. Monotonicity supplies its principal-stratum interpretation.

\begin{assumptionv}{ass:monotonicity}[Monotonicity of treatment receipt]
The potential receipts \(D(1)\) and \(D(0)\), under unit and zero encouragement respectively, satisfy
\[
D(1)\ge D(0)
\]
almost surely under \(P_n^F\), the primitive full-data law at sample-size index \(n\).
\end{assumptionv}

The no-defiers monotonicity condition makes individual receipt weakly increasing in encouragement \citep{ChenHuang2025Generalizing}. Under these instrument restrictions, the source receipt contrast represents a conditional complier share, and the source outcome contrast represents a conditional complier-weighted effect; their transported interpretation is given in \cref{obj:lem:transported-cace-identification}.

Transport is imposed through two equalities at target-almost-every covariate value. The first carries the complier-weighted outcome contrast across populations.

\begin{assumptionv}{ass:transport-complier-outcome}[Transport of complier outcomes]
Let \(C\) be the indicator of the complier principal stratum, let \(Y(d)\) be the potential outcome under receipt \(d\), and let \(P_n^F\) be the primitive full-data law at sample-size index \(n\). For \(P_{\mathrm T}\)-almost every covariate value \(x\), where \(P_{\mathrm T}\) is the law of one observed target covariate,
\[
\mathbb E_{P_n^F}[(Y(1)-Y(0))C\mid X=x,S=0]
=
\mathbb E_{P_n^F}[(Y(1)-Y(0))C\mid X=x,S=1],
\]
where \(S=0\) denotes target membership and \(S=1\) denotes source membership.
\end{assumptionv}

After aligning encouragement potential outcomes with receipt-mediated outcomes under exclusion and monotonicity, this complier-weighted outcome restriction corresponds to Assumption 2 of \citet{ChenHuang2025Generalizing}. It transports the conditional contribution to the effect numerator, including local complier prevalence. The second equality transports that prevalence itself.

\begin{assumptionv}{ass:transport-complier-share}[Transport of complier shares]
Let \(C\) be the indicator of the complier principal stratum, and let \(P_n^F\) be the primitive full-data law at sample-size index \(n\). For \(P_{\mathrm T}\)-almost every covariate value \(x\), where \(P_{\mathrm T}\) is the law of one observed target covariate,
\[
\mathbb E_{P_n^F}[C\mid X=x,S=0]
=
\mathbb E_{P_n^F}[C\mid X=x,S=1],
\]
where \(S=0\) denotes target membership and \(S=1\) denotes source membership.
\end{assumptionv}

The pointwise complier-share equality is the transport condition chosen for the present model and implies the averaged equality in Assumption 4 of \citet{ChenHuang2025Generalizing}. It is substantively plausible when source participation changes covariate composition but, within the measured covariate strata, does not change how encouragement shifts receipt. Together, the two equalities carry the numerator and denominator of the target complier effect through the target covariate distribution. The contribution developed below concerns estimation over the rough nuisance class and the honest expected-length frontier under these identification conditions.



The source contrast \leanref{sym:\Delta_A}{\ensuremath{\Delta_A}} and target-averaged contrast \leanref{sym:T_A}{\ensuremath{T_A}} separate the source conditional response from target covariate weighting. The target complier share \leanref{sym:\mu}{\ensuremath{\mu}} determines the normalization of the target complier effect \leanref{sym:\theta}{\ensuremath{\theta}}. Because both potential outcomes lie in \([0,1]\), the effect domain \leanref{sym:\Theta}{\ensuremath{\Theta}} is \([-1,1]\). Positive compliance gives the conditional target effect a well-defined denominator.

\begin{assumptionv}{ass:positive-strength}[Positive transported first stage]
The transported compliance share and first stage, \(\mu\), satisfies
\[
\mu>0.
\]
\end{assumptionv}

Target instrument relevance requires a positive transported complier share \citep{ChenHuang2025Generalizing}. The model permits this share to approach zero as the sample-size index varies. Under the preceding causal and transport restrictions, \cref{obj:lem:transported-cace-identification} gives
\[
\mu=T_D,
\qquad
\theta=\frac{T_Y}{T_D}.
\]
Transported compliance therefore determines both the population represented by the causal effect and the first stage that scales the observed outcome contrast. The bounded domain gives confidence intervals a common length scale even when that first stage is small.

\subsection{Global honesty and expected length}

We collect the sampling, regularity and causal restrictions into a single law-level class before defining coverage.

\begin{definitionv}{def:model-class}[Model class $\mathcal M_n$]
The model class \(\mathcal M_n\) consists of law triples
\[
P=(P_n^F,P_{\mathrm S},P_{\mathrm T}),
\]
where \(P_n^F\) is the primitive full-data law, \(P_{\mathrm S}\) is the law of one observed source record, and \(P_{\mathrm T}\) is the law of one observed target covariate, satisfying all of the following restrictions at sample-size index \(n\):
\begin{itemize}
\item \textbf{(Sampling.)} \cref{obj:ass:source-iid,obj:ass:target-iid,obj:ass:sample-independence}.
\item \textbf{(Covariate densities.)} \cref{obj:ass:source-density-bounds,obj:ass:target-density-bounds,obj:ass:source-density-holder,obj:ass:target-density-holder}.
\item \textbf{(Propensity and arm means.)} \cref{obj:ass:propensity-overlap,obj:ass:propensity-holder,obj:ass:arm-means-holder}.
\item \textbf{(Causal restrictions.)} \cref{obj:ass:receipt-consistency,obj:ass:outcome-consistency,obj:ass:instrument-randomization,obj:ass:exclusion,obj:ass:monotonicity}.
\item \textbf{(Transport and strength.)} \cref{obj:ass:transport-complier-outcome,obj:ass:transport-complier-share,obj:ass:positive-strength}.
\end{itemize}
\end{definitionv}

A law triple \leanref{sym:P}{\ensuremath{P}} in the model class \leanref{sym:\mathcal M_n}{\ensuremath{\mathcal M_n}} specifies both the causal target and the observed two-sample experiment. For probabilities and expectations involving a data-dependent interval, the observed law is \(P_{\mathrm S}^{\otimes n}\otimes P_{\mathrm T}^{\otimes n}\). The following conventions fix the coverage level, the estimator's sample-size threshold and the measure used to evaluate interval length.



The noncoverage probability \leanref{sym:\alpha}{\ensuremath{\alpha}} is fixed across the model class. The sample-size threshold \leanref{sym:n_0}{\ensuremath{n_0}} governs use of the cubic estimator, and restricted Lebesgue length \leanref{sym:\lambda_{\Theta}}{\ensuremath{\lambda_\Theta}} measures uncertainty on the causal domain. A generic random interval \leanref{sym:C_n}{\ensuremath{C_n}} is a measurable function of the two observed samples: formally, its graph \(\{(\omega,t):t\in C_n(\omega)\}\) is measurable in the product of the observed sample space and the Borel \(\sigma\)-field on \(\Theta\). We require its coverage to hold uniformly over the full model class.

\begin{definitionv}{def:honest-intervals}[Honest interval class \(\mathfrak H_n\)]
The class of uniformly honest connected intervals is
\[
\mathfrak H_n=
\left\{
C_n:
C_n\subseteq\Theta\ \text{is a measurable connected interval},
\quad
\inf_{P\in\mathcal M_n}P(\theta\in C_n)\ge 1-\alpha
\right\}.
\]
Here \(\Theta\) is the bounded domain of the target-population complier effect \(\theta\), \(\mathcal M_n\) is the model class, and \(\alpha\) is the prescribed noncoverage probability.
\end{definitionv}

The honest interval class \leanref{sym:\mathfrak H_n}{\ensuremath{\mathfrak H_n}} imposes finite-sample coverage at every law in \(\mathcal M_n\), including laws with arbitrarily small positive transported compliance. Connectedness makes the reported set an interval, and containment in \(\Theta\) bounds its length by two. To evaluate how its expected length changes with first-stage strength, we next select a strength slice.

\begin{definitionv}{def:strength-slice}[First-stage strength slice]
The first-stage strength slice, with \(\mu(P)\) denoting the transported first stage under \(P\), is defined for real parameters \(c_f,C_f,L,a\) and every \(n\in\mathbb N\) by
\[
\mathcal M_n(a)
=
\left\{
P\in\mathcal M_n:
0<a\le\frac14,\quad
a\le\mu(P)\le\frac14
\right\}.
\]
Here \(\mathcal M_n\) is the model class of \cref{obj:def:model-class} with parameters \(c_f,C_f,L\).
\end{definitionv}

The evaluation floor \leanref{sym:a}{\ensuremath{a}} indexes the strength slice \leanref{sym:\mathcal M_n(a)}{\ensuremath{\mathcal M_n(a)}}, whose transported first stages lie between \(a\) and \(1/4\). It enters the laws over which expected length is evaluated. Coverage continues to be required over all of \(\mathcal M_n\), and the attaining interval in \cref{obj:def:score-interval} uses a single construction across these slices. This ordering separates global validity from the strength-specific evaluation of precision.

We express the resulting optimal expected-length criterion as follows.

\begin{definitionv}{def:length-frontier}[Minimax length frontier \(L_n(a)\)]
The minimax expected-length frontier is
\[
L_n(a)
=
\inf_{C_n\in\mathfrak H_n}
\sup_{P\in\mathcal M_n(a)}
\mathbb E_P\!\left[\lambda_{\Theta}(C_n)\right],
\qquad n\ge n_0,\quad 0<a\le\frac14,
\]
where \(\mathfrak H_n\) is the class of uniformly honest connected intervals, \(\mathcal M_n(a)\) is the model slice with transported first stage between \(a\) and \(1/4\), and \(\lambda_{\Theta}\) denotes Lebesgue length restricted to the effect domain \(\Theta\).
\end{definitionv}

The minimax expected length \leanref{sym:L_n(a)}{\ensuremath{L_n(a)}} compares globally honest intervals by their worst expected length on a specified strength slice. Its infimum ranges over procedures satisfying the full coverage requirement in \cref{obj:def:honest-intervals}, while its supremum evaluates precision on \cref{obj:def:strength-slice}. This criterion links the roughness of the observed nuisance functions to the uncertainty remaining about a transported complier effect at a given first-stage scale.
